\documentclass{zhvt-classic}

% 使用本地源流明体（GenRyuMin2 TW）按文件路径加载，不依赖系统字体名
% 版框（内框／界栏／书口侧栏／鱼尾）改用与版心印章同一印泥棕色：
%   印章图片 yshs.png 实色 = RGB(135,68,52) = HTML 874434
\usepackage{xcolor}
\usepackage{tikz}
% sealbrown 双注册：xcolor 供 TikZ/pgf，L3 供 cls 版框
\definecolor{sealbrown}{HTML}{874434}
\ExplSyntaxOn
\color_set:nnn { sealbrown } { HTML } { 874434 }
\ExplSyntaxOff

\zhvtset{
  grid color       = sealbrown,
  main font        = GenRyuMin2TW-SB,
  main font path   = C:/Users/Administrator/texmf/fonts/opentype/public/GenRyuMin2TW/,
jiazhu font      = GenRyuMin2TW-M,
  jiazhu font path = C:/Users/Administrator/texmf/fonts/opentype/public/GenRyuMin2TW/,
  % 页码框容纳 4 字：正常页「九十七」+前置空 1 格=4 行恰好；
  % 100+ 页「一百零一」共 5 行超出，末字「一/二/三」被框裁掉（2026-10-04 实测），故放宽到 5
  shukou page number chars = 5,
  % 书口书名框高从默认 6 字减到 4 字，把省出的 2 字高拨给 chapter 框，
  % 使 6 字书名+卷名放进 9 字高的 chapter 框后不再压线（2026-10-05）
  shukou title chars = 4,
}

\title{史记}[刘敬叔孙通列传]
\author{汉·司马迁}
\maker{闲叟漫记·重製}


% ===== 索隐述赞样式（2026-11-08 四次定稿 · 愚大人亲手改定参数）=====
% ⚠ 竖排前提：自定义字体家族必须带 RawFeature={vertical:+vert}；
%   simhei.ttf 加 +vhal 会整体加载失败静默回退，只写 +vert；simfang 同理。
\newCJKfontfamily{\shuzanfont}[Path=C:/Windows/Fonts/,RawFeature={vertical:+vert}]{simhei.ttf}
\newCJKfontfamily{\shuzanbody}[Path=C:/Windows/Fonts/,RawFeature={vertical:+vert}]{simfang.ttf}
% 题头：152x30pt 圆角 6pt 棕底白字；baseline (0,2) 使整框落入界行格子不压版框线；
% 白字节点 x 不补偿、y +8.2pt（像素级墨迹实测 dx/dy≈0，L/R 留白约 20pt）。
\newcommand{\shuzanhead}[1]{%
  \begin{tikzpicture}[baseline={(0,2)}]
    \fill[sealbrown, rounded corners=6pt]
      (-76pt,-15pt) rectangle (76pt,15pt);
    \node[text=white, inner sep=0pt, anchor=center] at (0pt,8.2pt)
      {\shuzanfont\fontsize{23}{23}\selectfont #1};
  \end{tikzpicture}%
}
% ★ 修正书口卷名位置：chapter 框起始位置改回 cls 默认值，避免全局影响。
%   只在每个卷首页通过 \juan 命令里的空命令微调该页位置（安全，不会跨页污染）。
\ExplSyntaxOn
\AtBeginDocument{
  \dim_set:Nn \l__zhvt_shukou_chapter_title_start_dim
    { \l__zhvt_shukou_title_height_dim + \l__zhvt_font_size_dim * 2 }
  \dim_set:Nn \l__zhvt_shukou_page_number_start_dim
    { \l__zhvt_shukou_chapter_title_start_dim + \l__zhvt_shukou_chapter_title_height_dim }
}
\ExplSyntaxOff

\begin{document}

% ★★ 版心（书口）排布覆写：把书名放到鱼尾之下（2026-10-04）
%
%   cls 默认版心自上而下：书名（title 静态框，在鱼尾之上）→ 鱼尾 → 篇名（chapter 动态框）→ 页码
%   本模板改为：          鱼尾（位置不动）→ 书名 → 篇名 → 页码分割线 → 页码
%
%   手法：把 cls 放在 title 静态框里的书名，改放进 chapter 动态框（鱼尾之下）；
%         title 静态框留空。书名取 \l__zhvt_title_str，即 \title{...} 的第一参数，
%         所以换书时只改 \title，不必在此重复写书名。
%
%   ⚠ 三条必须遵守的写法（2026-10-04 实测）：
%     ① 汉字必须用 \hbox{...} 包住。flowfram 会把框内容装进一个宽 14.4mm（= 一个 line_sep
%        = 书口栏宽）的纵向 minipage 并按该宽度折行，而整页靠 PDF /Rotate 90 显示、
%        框内文字本身不变向 —— 不包 \hbox 时「围炉夜话」4 字会被逐字折成 4 行，
%        在版面上就成了横排 4 字并横向溢出版框。
%     ② 竖排下**一个 \hbox 就是一整列**。书名与篇名要同处一列，必须放进同一个 \hbox；
%        分成两个 \hbox 时，后一个会被当成新的一列、甩到隔壁栏去
%        （实测「序」跑到书口栏外侧 14.4mm 的正文列上）。中间的 \quad 在框内是横向空白，
%        显示出来正好是书名与篇名之间的纵向间隔。
%     ③ ★ 内容盒要写成 **与框等宽的定宽盒** `\hbox to \l__zhvt_shukou_chapter_title_height_dim{…\hfil}`。
%        实测：chapter 框是把内容盒沿纵向**居中**摆放的，内容宽度一变，书名就跟着漂：
%        自然宽度的 \hbox（书名 38.4mm）落在 x=157.2mm，加篇名后（57.6mm）落在 x=152.4mm，
%        目录页（无篇名）与正文页的书名就差半格。做成等宽盒后，无论有没有篇名，
%        内容固定从框顶起排，实测两页均落在 x=142.8mm，不再漂。
\ExplSyntaxOn
% 把书口 chapter 框（书名+篇名）起始位置下移 1 字高，避免 4 字卷名顶到鱼尾
% ⚠ cls 的 \__zhvt_calculate_dimensions: 在 \AtBeginDocument 里执行，会重算该 dim；
%   必须把覆写挂到 begin document 之后，否则会被二次覆盖（2026-10-05 实测）。
% (旧方案已移到导言区，直接修改 chapter 框起始位置)
% 卷标题命令 #1=书口篇名(鱼尾) #2=正文卷名（含卷序，卷名内空格用 ~）
%
% ★★ 目录页两项修订（2026-10-05，愚大人指出）★★
%   ① 页码起点统一对齐：
%      原状 —— 目录项写成 \par\hspace{4em}卷名\hspace{4em}页码，页码起点 = 卷名实宽 + 4em。
%      卷名实宽随字数变：「卷一~梁惠王上」6 字 +~、「卷七~离娄上」5 字 +~，
%      差 1 个字位 → 卷七~卷十 4 条的页码整体比其余 10 条高 27.21pt（一个字），看着上下不齐。
%      改法 —— 卷名补全角空格（宽 = 1 字 = 27.215pt）补足到 **7 字符位**再写进目录，
%      14 条页码起点遂全部落在同一条竖线上（起点对齐；末字随页码字数自然参差）。
%      ⚠ 只对目录里的副本 \l_tmpc_tl 补，正文卷名仍用 \l_tmpa_tl 原样 —— 正文像素不受影响。
%   ② 卷名整体上移半字位：前导 \hspace 由 4em 改 3.5em。
%      实测原状目录块在版框内偏下（天头留白 139.6pt / 地脚留白 110.7pt），
%      上移 0.5em = 13.6pt 后上下基本均分，视觉居中。
\NewDocumentCommand { \juan } { mm } {
	\clearpage
	\tl_set:Nn \l__zhvt_chapter_title_tl { #1 }
	\tl_set:Nn \l_tmpa_tl { #2 }
	\tl_set_eq:NN \l_tmpc_tl \l_tmpa_tl
	\int_step_inline:nn { 7 - \str_count:n { #2 } }
	  { \tl_put_right:Nn \l_tmpc_tl { 　 } }
	\tl_set:Nx \l_tmpb_tl { \zhnum { page } }
	\tl_put_left:Nn \l_tmpb_tl { \hspace { 4em } }
	\tl_put_left:Nn \l_tmpb_tl { \l_tmpc_tl }
	\tl_put_left:Nn \l_tmpb_tl { \par \hspace { 3.5em } }
	\addtocontents { toc } { \l_tmpb_tl }
	\skip_horizontal:n { \l__zhvt_font_size_dim * 4 }
	\l_tmpa_tl
	\par
	\thispagestyle { empty }
}
\cs_set_protected:Npn \__zhvt_set_shukou_content: {
	% ★ 鱼尾上方留白（title 静态框）嵌入闲叟钤印：
	%   整页 /Rotate 90，故此处把图用 \rotatebox{90} 转正为竖章；
	%   尺寸卡在版心条幅内（书内约 7mm 宽 × 33.8mm 高，留余量不压鱼尾）。
	\setstaticcontents * { title } {
		\vbox_to_ht:nn { \l__zhvt_baseline_skip_dim } {
			\vskip 14.2pt
			\hbox_to_wd:nn { \l__zhvt_shukou_title_height_dim } {
				\hfil
				\rotatebox [ origin = c ] { 90 } { \includegraphics [ width = 7mm , height = 33.8mm ] { yshs.png } }
				\hfil
			}
		}
	}
% ★ 版心（书口）统一改法（2026-11-16 愚大人定）：18pt＋框内居中＋书口文字同版框色
%   ① 原 27.2pt 时 7 字长篇名（如「廉颇蔺相如列传」）末字溢出 chapter 框、压住页码；
%      改 \fontsize{18pt}{18pt}，前后 \hfil 令字串在框内居中（不随篇名字数漂移）。
%   ② \vbox{\vskip 8.2pt} 定量下移，抵消缩字号后栏宽方向 8.2pt 上偏（像素实测 y 中心 543.3 = 原版）。
%      ⚠ \raisebox / \vspace* 在 flowfram 框内无效；换字号须同步换补偿：16pt→10.1、19pt→7.4。
%   ③ chapter 与 page~number 均置 \color{sealbrown}（HTML 874434），与版框／界行同色；正文仍黑。
%   ⚠ 页码字位、钤印、鱼尾、界行、正文一律未动；五部合订本同此口径。
	\setdynamiccontents * { chapter }
	{ \color{sealbrown} \vbox { \vskip 8.2pt \hbox to \l__zhvt_shukou_chapter_title_height_dim
		{ \hfil \fontsize{18pt}{18pt}\selectfont \l__zhvt_title_str \quad \l__zhvt_chapter_title_tl \hfil } } }
	\setdynamiccontents * { page~number } { \color{sealbrown} \quad \zhnum { page } }
}
\__zhvt_set_shukou_content:
\ExplSyntaxOff

\maketitle

% 单篇，无多卷目录，不渲染目录页

\juan{刘敬叔孙通列传}{卷九十九~刘敬叔孙通列传第三十九}
\hspace*{2em}刘敬者。齐人也。汉五年。戍陇西。过洛阳。高帝在焉。娄敬脱輓辂。衣其羊裘。见齐人虞将军曰。﹁臣原见上言便事。﹂虞将军欲与之鲜衣。娄敬曰。﹁臣衣帛。衣帛见。衣褐。衣褐见。终不敢易衣。﹂於是虞将军入言上。上召入见。赐食。\\
\hspace*{2em}已而问娄敬。娄敬说曰。﹁陛下都洛阳。岂欲与周室比隆哉。﹂上曰。﹁然。﹂娄敬曰。﹁陛下取天下与周室异。周之先自后稷。尧封之邰。积德累善十有馀世。公刘避桀居豳。太王以狄伐故。去豳。杖马箠居岐。国人争随之。及文王为西伯。断虞芮之讼。始受命。吕望。伯夷自海滨来归之。武王伐纣。不期而会孟津之上八百诸侯。皆曰纣可伐矣。遂灭殷。成王即位。周公之属傅相焉。乃营成周洛邑。以此为天下之中也。诸侯四方纳贡职。道里均矣。有德则易以王。无德则易以亡。凡居此者。欲令周务以德致人。不欲依阻险。令後世骄奢以虐民也。及周之盛时。天下和洽。四夷乡风。慕义怀德。附离而并事天子。不屯一卒。不战一士。八夷大国之民莫不宾服。效其贡职。及周之衰也。分而为两。天下莫朝。周不能制也。非其德薄也。而形势弱也。今陛下起丰沛。收卒三千人。以之径往而卷蜀汉。定三秦。与项羽战荥阳。争成皋之口。大战七十。小战四十。使天下之民肝脑涂地。父子暴骨中野。不可胜数。哭泣之声未绝。伤痍者未起。而欲比隆於成康之时。臣窃以为不侔也。且夫秦地被山带河。四塞以为固。卒然有急。百万之众可具也。因秦之故。资甚美膏腴之地。此所谓天府者也。陛下入关而都之。山东虽乱。秦之故地可全而有也。夫与人斗。不搤其亢。拊其背。未能全其胜也。今陛下入关而都。案秦之故地。此亦搤天下之亢而拊其背也。﹂\\
\hspace*{2em}高帝问群臣。群臣皆山东人。争言周王数百年。秦二世即亡。不如都周。上疑未能决。及留侯明言入关便。即日车驾西都关中。\\
\hspace*{2em}於是上曰。﹁本言都秦地者娄敬。﹃娄﹄者乃﹃刘﹄也。﹂赐姓刘氏。拜为郎中。号为奉春君。\\
\hspace*{2em}汉七年。韩王信反。高帝自往击之。至晋阳。闻信与匈奴欲共击汉。上大怒。使人使匈奴。匈奴匿其壮士肥牛马。但见老弱及羸畜。使者十辈来。皆言匈奴可击。上使刘敬复往使匈奴。还报曰。﹁两国相击。此宜夸矜见所长。今臣往。徒见羸瘠老弱。此必欲见短。伏奇兵以争利。愚以为匈奴不可击也。﹂是时汉兵已逾句注。二十馀万兵已业行。上怒。骂刘敬曰。﹁齐虏。以口舌得官。今乃妄言沮吾军。﹂械系敬广武。遂往。至平城。匈奴果出奇兵围高帝白登。七日然後得解。高帝至广武。赦敬。曰。﹁吾不用公言。以困平城。吾皆已斩前使十辈言可击者矣。﹂乃封敬二千户。为关内侯。号为建信侯。\\
\hspace*{2em}高帝罢平城归。韩王信亡入胡。当是时。冒顿为单于。兵彊。控弦三十万。数苦北边。上患之。问刘敬。刘敬曰。﹁天下初定。士卒罢於兵。未可以武服也。冒顿杀父代立。妻群母。以力为威。未可以仁义说也。独可以计久远子孙为臣耳。然恐陛下不能为。﹂上曰。﹁诚可。何为不能。顾为柰何。﹂刘敬对曰。﹁陛下诚能以適长公主妻之。厚奉遗之。彼知汉適女送厚。蛮夷必慕以为阏氏。生子必为太子。代单于。何者。贪汉重币。陛下以岁时汉所馀彼所鲜数问遗。因使辩士风谕以礼节。冒顿在。固为子婿。死。则外孙为单于。岂尝闻外孙敢与大父抗礼者哉。兵可无战以渐臣也。若陛下不能遣长公主。而令宗室及後宫诈称公主。彼亦知。不肯贵近。无益也。﹂高帝曰。﹁善。﹂欲遣长公主。吕后日夜泣。曰。﹁妾唯太子。一女。柰何弃之匈奴。﹂上竟不能遣长公主。而取家人子名为长公主。妻单于。使刘敬往结和亲约。\\
\hspace*{2em}刘敬从匈奴来。因言﹁匈奴河南白羊。楼烦王。去长安近者七百里。轻骑一日一夜可以至秦中。秦中新破。少民。地肥饶。可益实。夫诸侯初起时。非齐诸田。楚昭。屈。景莫能兴。今陛下虽都关中。实少人。北近胡寇。东有六国之族。宗彊。一日有变。陛下亦未得高枕而卧也。臣原陛下徙齐诸田。楚昭。屈。景。燕。赵。韩。魏後。及豪桀名家居关中。无事。可以备胡。诸侯有变。亦足率以东伐。此彊本弱末之术也﹂。上曰。﹁善。﹂乃使刘敬徙所言关中十馀万口。\\
\hspace*{2em}叔孙通者。薛人也。秦时以文学徵。待诏博士。数岁。陈胜起山东。使者以闻。二世召博士诸儒生问曰。﹁楚戍卒攻蕲入陈。於公如何。﹂博士诸生三十馀人前曰。﹁人臣无将。将即反。罪死无赦。原陛下急发兵击之。﹂二世怒。作色。叔孙通前曰。﹁诸生言皆非也。夫天下合为一家。毁郡县城。铄其兵。示天下不复用。且明主在其上。法令具於下。使人人奉职。四方辐辏。安敢有反者。此特群盗鼠窃狗盗耳。何足置之齿牙间。郡守尉今捕论。何足忧。﹂二世喜曰。﹁善。﹂尽问诸生。诸生或言反。或言盗。於是二世令御史案诸生言反者下吏。非所宜言。诸言盗者皆罢之。乃赐叔孙通帛二十匹。衣一袭。拜为博士。叔孙通已出宫。反舍。诸生曰。﹁先生何言之谀也。﹂通曰。﹁公不知也。我几不脱於虎口。﹂乃亡去。之薛。薛已降楚矣。及项梁之薛。叔孙通从之。败於定陶。从怀王。怀王为义帝。徙长沙。叔孙通留事项王。汉二年。汉王从五诸侯入彭城。叔孙通降汉王。汉王败而西。因竟从汉。\\
\hspace*{2em}叔孙通儒服。汉王憎之。乃变其服。服短衣。楚制。汉王喜。\\
\hspace*{2em}叔孙通之降汉。从儒生弟子百馀人。然通无所言进。专言诸故群盗壮士进之。弟子皆窃骂曰。﹁事先生数岁。幸得从降汉。今不能进臣等。专言大猾。何也。﹂叔孙通闻之。乃谓曰。﹁汉王方蒙矢石争天下。诸生宁能斗乎。故先言斩将搴旗之士。诸生且待我。我不忘矣。﹂汉王拜叔孙通为博士。号稷嗣君。\\
\hspace*{2em}汉五年。已并天下。诸侯共尊汉王为皇帝於定陶。叔孙通就其仪号。高帝悉去秦苛仪法。为简易。群臣饮酒争功。醉或妄呼。拔剑击柱。高帝患之。叔孙通知上益厌之也。说上曰。﹁夫儒者难与进取。可与守成。臣原徵鲁诸生。与臣弟子共起朝仪。﹂高帝曰。﹁得无难乎。﹂叔孙通曰。﹁五帝异乐。三王不同礼。礼者。因时世人情为之节文者也。故夏。殷。周之礼所因损益可知者。谓不相复也。臣原颇采古礼与秦仪杂就之。﹂上曰。﹁可试为之。令易知。度吾所能行为之。﹂\\
\hspace*{2em}於是叔孙通使徵鲁诸生三十馀人。鲁有两生不肯行。曰。﹁公所事者且十主。皆面谀以得亲贵。今天下初定。死者未葬。伤者未起。又欲起礼乐。礼乐所由起。积德百年而後可兴也。吾不忍为公所为。公所为不合古。吾不行。公往矣。无汙我。﹂叔孙通笑曰。﹁若真鄙儒也。不知时变。﹂\\
\hspace*{2em}遂与所徵三十人西。及上左右为学者与其弟子百馀人为绵蕞野外。习之月馀。叔孙通曰。﹁上可试观。﹂上既观。使行礼。曰。﹁吾能为此。﹂乃令群臣习肄。会十月。\\
\hspace*{2em}汉七年。长乐宫成。诸侯群臣皆朝十月。仪。先平明。谒者治礼。引以次入殿门。廷中陈车骑步卒卫宫。设兵张旗志。传言﹁趋﹂。殿下郎中侠陛。陛数百人。功臣列侯诸将军军吏以次陈西方。东乡。文官丞相以下陈东方。西乡。大行设九宾。胪传。於是皇帝辇出房。百官执职传警。引诸侯王以下至吏六百石以次奉贺。自诸侯王以下莫不振恐肃敬。至礼毕。复置法酒。诸侍坐殿上皆伏抑首。以尊卑次起上寿。觞九行。谒者言﹁罢酒﹂。御史执法举不如仪者辄引去。竟朝置酒。无敢讙譁失礼者。於是高帝曰。﹁吾乃今日知为皇帝之贵也。﹂乃拜叔孙通为太常。赐金五百斤。\\
\hspace*{2em}叔孙通因进曰。﹁诸弟子儒生随臣久矣。与臣共为仪。原陛下官之。﹂高帝悉以为郎。叔孙通出。皆以五百斤金赐诸生。诸生乃皆喜曰。﹁叔孙生诚圣人也。知当世之要务。﹂\\
\hspace*{2em}汉九年。高帝徙叔孙通为太子太傅。汉十二年。高祖欲以赵王如意易太子。叔孙通谏上曰。﹁昔者晋献公以骊姬之故废太子。立奚齐。晋国乱者数十年。为天下笑。秦以不蚤定扶苏。令赵高得以诈立胡亥。自使灭祀。此陛下所亲见。今太子仁孝。天下皆闻之。吕后与陛下攻苦食啖。其可背哉。陛下必欲废適而立少。臣原先伏诛。以颈血汙地。﹂高帝曰。﹁公罢矣。吾直戏耳。﹂叔孙通曰。﹁太子天下本。本一摇天下振动。柰何以天下为戏。﹂高帝曰。﹁吾听公言。﹂及上置酒。见留侯所招客从太子入见。上乃遂无易太子志矣。\\
\hspace*{2em}高帝崩。孝惠即位。乃谓叔孙生曰。﹁先帝园陵寝庙。群臣莫习。﹂徙为太常。定宗庙仪法。及稍定汉诸仪法。皆叔孙生为太常所论箸也。\\
\hspace*{2em}孝惠帝为东朝长乐宫。及间往。数跸烦人。乃作衤复道。方筑武库南。叔孙生奏事。因请间曰。﹁陛下何自筑衤复道高寝。衣冠月出游高庙。高庙。汉太祖。柰何令後世子孙乘宗庙道上行哉。﹂孝惠帝大惧。曰。﹁急坏之。﹂叔孙生曰。﹁人主无过举。今已作。百姓皆知之。今坏此。则示有过举。原陛下原庙渭北。衣冠月出游之。益广多宗庙。大孝之本也。﹂上乃诏有司立原庙。原庙起。以衤复道故。\\
\hspace*{2em}孝惠帝曾春出游离宫。叔孙生曰。﹁古者有春尝果。方今樱桃孰。可献。原陛下出。因取樱桃献宗庙。﹂上乃许之。诸果献由此兴。\\
\hspace*{2em}太史公曰。语曰﹁千金之裘。非一狐之腋也。台榭之榱。非一木之枝也。三代之际。非一士之智也﹂。信哉。夫高祖起微细。定海内。谋计用兵。可谓尽之矣。然而刘敬脱輓辂一说。建万世之安。智岂可专邪。叔孙通希世度务。制礼进退。与时变化。卒为汉家儒宗。﹁大直若诎。道固委蛇﹂。盖谓是乎。\\

% ===== 索隐述赞（司马贞《史记索隐》）· 统一样式 =====
% 题头「索隐述赞曰」五字并入反白框；正文每四句一列，定宽盒 \hfil 居列中；
% \fontsize{24}{41} 列距对齐全书界行（121+41k）；\clearpage 使述赞独占一页。
\clearpage
\shuzanhead{索隐述赞曰}\linebreak
{\shuzanbody\fontsize{24}{41}\selectfont
\hbox to \hsize{\hfil 厦藉众幹。裘非一狐。委辂献说。釂蕝陈书。\hfil}\linebreak
\hbox to \hsize{\hfil 皇帝始贵。车驾西都。既安太子。又和匈奴。\hfil}\linebreak
\hbox to \hsize{\hfil 奉春。稷嗣。其功可图。\hfil}%
\par}
\end{document}
